\documentclass[12pt]{article}
\usepackage[utf8]{inputenc} % Sin el guion
\usepackage[T1]{fontenc}    % Ayuda a que las tildes se copien bien del PDF
\usepackage[spanish]{babel}
\usepackage{amsmath}
\usepackage{amssymb}
\usepackage{geometry}
\usepackage[backend=biber,style=numeric]{biblatex}
\usepackage{csquotes}
% \bibliography{referencias}
\geometry{margin=1in}
\usepackage{listings}
\usepackage{xcolor}

\usepackage{graphicx}
\usepackage{float}      % opcional, para usar [H]
\usepackage{hyperref}   % opcional, para referencias clickeables

\lstset{
    language=Python,
    basicstyle=\ttfamily\footnotesize,
    keywordstyle=\color{blue},
    commentstyle=\color{gray},
    stringstyle=\color{black},
    numbers=left,
    numberstyle=\tiny,
    frame=single,
    rulecolor=\color{black},
    breaklines=true,
    showstringspaces=false
}

\title{Actividad 4. Modelos de prediccion y Metricas de Error}
\author{Ing. Isaias Garcia Zanella}
\date{\today}

\begin{document}

\maketitle

\section{Introducción}
\section{Marco Teórico}

\subsection{Metricas de Error}
Las metricas de evaluacion para datos continuos resultan utiles para conocer si los valores aproximados o predichos por un modelo se acercan o no a los valores reales. Esta diferencia entre los valores reales y lor predichos es el error.
Cada metrica que se comentara a continuacion trata este error de manera singular, aplicando diversas ecuaciones que son convenientes para diversos modelos.A continuacion mencionaremos algunas metricas que pueden utilizarse:
\begin{enumerate}
    \item Error medio Cuadratico (MSE)
    \item Error Medio Absoluto (MAE)
    \item Error Relativo Cuadratico (RSE)
    \item Error Relativo Absoluto (RAE)
    \item Coeficiente de Correlacion (CC)
\end{enumerate}

\subsubsection{ Mean Squared Error (MSE y RMSE)}
MSE es comunmente el metodo mas utilizafo para poder predecir el error. Tambien se utiliza su raiz cuadrada, llamada <<Root Mean-Squared Error>>-RMSE-
Si consideramos el valor real (actual) como $a$ y el valor predicho como $p$, MSE y RMSE pueden calcularse como la diferencia entre el valor actual y el predicho al cuadrado entre el numero de valores (n) de la siguiente manera:
\begin{equation}
    MSE = \frac{1}{n} \sum_{}^{} (i=1)^n (y_i - \hat{y_i})^2
\end{equation}
\begin{equation}
    RMSE = \sqrt{\frac{1}{n} \sum_{}^{}(i=1)^n(y_i - \hat{y_i})^2}
\end{equation}

Donde:
\begin{itemize}
    \item $y_i$ = valor real
    \item $\hat{y_i}$ = valor predicho
    \item $n$ = numero de observaciones
\end{itemize}

El error medio cuadratico permite evaluar el rendimiento de multiples modelos en un problema de prediccion cuando se trata de datos continuos. MSE varia de un rango de $[0,\infty)$ siendo
un menor valor de MSE, un mejor rendimiento del modelo. Este metodo penaliza mas los errores grandes debido al cuadrado  y es ampliamente utilizado y bien entendido. Sin embargo , tiene algunas limitaciones, entre ellas, que MSE es sensible a valores atipicos << outliers>> y tampoco proporciona informacion sobre direccion del error.


\subsubsection{Mean Aboslute Error(MAE)}
Es una medida que solo promedia las magnitudes de cada error de manera individual sin tener en cuenta su signo. Tiene una ventaja con respecto a MSE; con MAE los outliers no son exagerados y se trata todos los errores de la misma manera de acuerdo a su magnitud.
El error medio absoluto se calcula de las siguiente manera:
\begin{equation}
    MAE = \frac{1}{n} \sum_{}^{}(i=1)^n |y_i - \hat{y_i}|
\end{equation}


\subsubsection{Relative squared error (RSE)}
Con el Error Relativo Cuadratico (RSE) se normaliza, lo cual permite una comparacion entre diferentes conjuntos de datos y es independiente de la escala de los datos. Las desventajas de este metodo es que es menos intituivo que MSE o MAE y requiere calcular la media de los datos reales, por lo que es un poco mas complicado de calcular y hacer sentido del resultado.
\begin{equation}
    MAE = \frac{\sum_{}^{}(i=1)^n (y_i - \hat{y_i})^2}{\sum_{}^{}(i=1)^n (y_i - \bar{y})^2} 
\end{equation}
Donde:
\begin{itemize}
    \item $y_i$ = valor real de la observacion 
    \item $\hat{y_i}$ = valor predicho de la observacion 
    \item $\bar{y}$ + media de los valores reales : $\bar{y} = (\frac{1}{n}) \sum_{(i=1)^n}{} y_i $
\end{itemize}

\subsubsection{Relative absolute error (RAE)}
El Error Relativo Absoluto (RAE) es el valor del error absoluto, normaliado de la misma manera que con MAE. RAE mide la fuerza de la relacion lineal
y; como es una especie de normalizacion, sus valores se encuentran entre -1 y 1 por lo que este metodo es de facil interpretacion. Tiene algunas desventajas, por ejemplo, solo mide correlacion lineal y no indica la magnitud del error.

\begin{equation}
    RAE = \frac{\sum_{i=1}^{n} | y_i - \hat{y_i} |}{\sum_{i=1}^{n} | y_i - \bar{y} |}
\end{equation}

\subsubsection{Correlation Coefficient (CC)}
El coeficiente de correlacion (CC) mide la correlacion estadistica los valores actuales y los valores predichos. El coeficiente de correlacion varia de 1 para un resultado de correlacion es perfecta pero negativa. Mide
la fuerza de la relacion lineal y debido a su rango de valores, es facil interpretar esta metrica. Como deventaja, no indica la magnitud del error y puede ser alto incluso con predicciones sesgadas. 

\begin{equation}
    CC = \frac{\sum_{}^{}[(y_i - \bar{y})(\hat{y_i} - \hat{\bar{y}})]}{\sqrt{[\sum_{}^{}(yi - \bar{y})^2 \times \sum_{}^{}(\hat{y}i - \bar{\hat{y}})^2]}}
\end{equation}

\section{Materiales y Métodos}
\section{Resultados y Discusión}
\section{Conclusión}
% \printbibliography[title={ Bibliográficas}]
\end{document}